\subsection{The coefficient image and its constant block}\label{app:compression}

We refine the preceding argument by counting the coefficients of
$Q_i$, rather than all the moments. The same disjointness-matrix
coefficient count is used in~\cite[Section 3]{Guo}. The additional
identity~\eqref{eq:constant-block} removes its constant block on
the critical locus. We give the details over an algebraically closed
field of characteristic zero or characteristic $p>\max(t,s)$.

For $\varnothing\ne S\subseteq[t]$ with $|S|\le d-1$, put
\[
 h_S=\sum_{\pi\vdash S}\mu(\pi)\prod_{B\in\pi}p_B,
 \qquad h_\varnothing=1.
\]
For $\varnothing\ne T\subseteq[t]$, $|T|\le d$, define
\[
 c_T=\sum_{\substack{S\subseteq[t]\setminus T\\|S|=d-|T|}}h_S.
\]
Collecting coefficients in~\eqref{eq:partitionQ} gives
\begin{equation}\label{eq:collected-Q}
 Q_i(u;p)=\sum_{\substack{J\subseteq[t]\setminus\{i\}\\|J|\le d-1}}
       (-1)^{|J|}|J|!\,c_{J\cup\{i\}}u_J.
\end{equation}
To see this, choose the blocks supplying a factor $-u_B$ and let
$J$ be their union. The partitions of the remaining rows contribute
$h_S$. On $J$, summing the products of cyclic-order counts over its
partitions gives $|J|!$, with sign $(-1)^{|J|}$. Summing over the
remaining $S$ proves the identity. The leading coefficients, for
$|J|=d-1$, have $c_{J\cup\{i\}}=1$.

For $|T|=j$, the vector $(c_T)$ is a linear image of $(h_S)_{|S|=d-j}$
under the disjointness matrix on $j$-subsets and $(d-j)$-subsets
of $[t]$. Its rank is
\[
 \min\left\{\binom tj,\binom t{d-j}\right\}
 =\binom t{\min(j,d-j)}.
\]
Here is a proof of the rank fact, including its characteristic range.
Transpose and complement subsets to reduce to the following assertion:
if $a\le v\le t-a$ and
\[
 \sum_{\substack{A\subseteq V\\|A|=a}}z_A=0
 \quad\hbox{for every $v$-subset $V\subseteq[t]$},
\]
then every $z_A=0$. Induct on $a$, the case $a=0$ being immediate.
For distinct $i,j$, subtract the equations for $U\cup\{i\}$ and
$U\cup\{j\}$, where $U$ ranges over the $(v-1)$-subsets avoiding
$i,j$. On the remaining $t-2$ elements the induction hypothesis
applies to the coefficients $z_{C\cup\{i\}}-z_{C\cup\{j\}}$:
$a-1\le v-1\le(t-2)-(a-1)$. They all vanish. The exchange graph
of $a$-subsets is connected, so every $z_A$ has one common value $z$.
The original equation now reads $\binom va z=0$; this binomial
coefficient is nonzero in the stated characteristics. This proves
the rank formula, a form of the classical inclusion-matrix rank
theorem also used in~\cite{Guo}.

The change from the independent moments $p_S$ to the $h_S$ is
triangular by $|S|$, with nonzero diagonal coefficient
$(-1)^{|S|-1}(|S|-1)!$. Thus the $h_S$ are independent coordinates.
Different $j$ use different degrees $d-j$, and their images occupy
different monomials in~\eqref{eq:collected-Q}. Consequently the
formal coefficient image has dimension exactly
\[
 R(t,d)=\sum_{j=1}^{d-1}\binom t{\min(j,d-j)}.
\]
The $j=1$ block consists of $t$ independent constant coefficients.

At an actual matrix, inclusion--exclusion identifies $h_S$ with the
sum of the matchings using exactly the row set $S$. For each $i$,
direct counting gives
\begin{equation}\label{eq:constant-block}
 \sum_{a=1}^s\partial_{ia}M_{t,s,d}=(s-d+1)c_{\{i\}}.
\end{equation}
A matching of size $d-1$ avoiding row $i$ is counted once for each
unused column. Since $s-d+1\ne0$ in the ground field, every
$c_{\{i\}}$ vanishes on the critical locus. Imposing this zero
block on the formal coefficient image leaves dimension
$R_{\rm crit}(t,d)=R(t,d)-t$.

Choose coordinates on each remaining disjointness image and use
them to reconstruct the coefficients of~\eqref{eq:collected-Q}.
Together with the $s(d-2)$ column elementary invariants used above,
these coordinates generate a polynomial parameter ring of dimension
$R_{\rm crit}(t,d)+s(d-2)$. Its universal column incidence still
contains every critical matrix with its actual coefficient data.
The leading part of $Q_i$ is unchanged, so the same degree-reducing
monic relations make its coordinate ring finite over this parameter
ring. Surjecting onto the critical-locus coordinate ring proves
\eqref{eq:critical-compressed}.

For use in the parameter estimates, the count has the explicit forms
\[
 \begin{split}
 R_{\rm crit}(t,2h)&=2\sum_{j=1}^{h-1}\binom tj+\binom th-t,\\
 R_{\rm crit}(t,2h+1)&=2\sum_{j=1}^{h}\binom tj-t.
 \end{split}
\]
The dimension calculation is exact for the formal coefficient image.
It does not assert that actual critical matrices fill that image,
or that the resulting upper bound on their dimension is sharp.
